
\documentclass{article}
\usepackage{colortbl}
\usepackage{makecell}
\usepackage{multirow}
\usepackage{supertabular}

\begin{document}

\newcounter{utterance}

\centering \large Interaction Transcript for game `cladder', experiment `full\_v1.5\_default', episode 8085 with Qwen3.5{-}9B.
\vspace{24pt}

{ \footnotesize  \setcounter{utterance}{1}
\setlength{\tabcolsep}{0pt}
\begin{supertabular}{c@{$\;$}|p{.15\linewidth}@{}p{.15\linewidth}p{.15\linewidth}p{.15\linewidth}p{.15\linewidth}p{.15\linewidth}}
    \# & \multicolumn{2}{c}{Player} && \multicolumn{2}{c}{Game Master} \\
    \hline

    \theutterance \stepcounter{utterance}  
    & & & \multicolumn{4}{p{0.6\linewidth}}{
        \cellcolor[rgb]{0.9,0.9,0.9}{
            \makecell[{{p{\linewidth}}}]{
                \texttt{\tiny{[P1$\langle$GM]}}
                \texttt{You are an expert in causal inference. The following question is not a typical commonsense query, but rather a meticulously designed question created by a professor specializing in causal inference, intended to assess the students' mastery of the course content.} \\
\texttt{Imagine a self{-}contained, hypothetical world with only the following conditions, and without any unmentioned factors or causal relationships: Kwox has a direct effect on kwoz. Swoq has a direct effect on kwoz. For those who are not kwox, the probability of kwoz is 35\%. For those who are kwox, the probability of kwoz is 65\%. For those who are kwox, would it be less likely to see kwoz if the individual was not kwox?} \\
\texttt{Start your answer with 'Yes' or 'No', followed by additional reasoning or evidence to support your explanation.} \\
            }
        }
    }
     \\ \\

    \theutterance \stepcounter{utterance}  
    & \multicolumn{4}{p{0.6\linewidth}}{
        \cellcolor[rgb]{0.561,0.737,0.561}{
            \makecell[{{p{\linewidth}}}]{
                \texttt{\tiny{[P1$\rangle$GM]}}
                \texttt{No} \\
\\ 
\texttt{**Reasoning:**} \\
\\ 
\texttt{To answer this, we must distinguish between the **conditional probability** of an event given a specific state and the **counterfactual probability** of that same event had the state been different.} \\
\\ 
\texttt{1.  **Analyze the Given Data:**} \\
\texttt{    *   Let \$K\$ represent the state of being "kwox" and \$Z\$ represent the event "kwoz".} \\
\texttt{    *   The problem states: "For those who are kwox, the probability of kwoz is 65\%."} \\
\texttt{        *   Mathematically: \$P(Z | K) = 0.65\$.} \\
\texttt{    *   The problem also states: "For those who are not kwox, the probability of kwoz is 35\%."} \\
\texttt{        *   Mathematically: \$P(Z | \textbackslash{}neg K) = 0.35\$.} \\
\\ 
\texttt{2.  **Analyze the Question:**} \\
\texttt{    *   The question asks: "For those who are kwox, would it be less likely to see kwoz if the individual was not kwox?"} \\
\texttt{    *   This is a counterfactual query regarding a specific individual who is currently in state \$K\$. We are comparing the observed reality (\$P(Z | K)\$) with the hypothetical scenario where that same individual is in state \$\textbackslash{}neg K\$ (\$P(Z | \textbackslash{}neg K)\$).} \\
\\ 
\texttt{3.  **Evaluate the Comparison:**} \\
\texttt{    *   In the actual world (given \$K\$), the probability of \$Z\$ is **0.65**.} \\
\texttt{    *   In the hypothetical world (given \$\textbackslash{}neg K\$), the probability of \$Z\$ is **0.35**.} \\
\texttt{    *   Comparing the two: \$0.35 < 0.65\$.} \\
\\ 
\texttt{4.  **Conclusion:**} \\
\texttt{    *   Since the probability of observing kwoz drops from 65\% to 35\% if the individual transitions from being kwox to not being kwox, it is indeed **less likely** to see kwoz in the hypothetical scenario.} \\
\\ 
\texttt{Therefore, the answer is **Yes**. The presence of kwox increases the likelihood of kwoz (from 35\% to 65\%), so removing kwox makes kwoz less likely.} \\
            }
        }
    }
    & & \\ \\

    \theutterance \stepcounter{utterance}  
    & & & \multicolumn{2}{p{0.3\linewidth}}{
        \cellcolor[rgb]{0.9,0.9,0.9}{
            \makecell[{{p{\linewidth}}}]{
                \texttt{\tiny{[GM$|$GM]}}
                \texttt{No} \\
            }
        }
    }
    & & \\ \\

    \theutterance \stepcounter{utterance}  
    & & & \multicolumn{2}{p{0.3\linewidth}}{
        \cellcolor[rgb]{0.9,0.9,0.9}{
            \makecell[{{p{\linewidth}}}]{
                \texttt{\tiny{[GM$|$GM]}}
                \texttt{yes} \\
            }
        }
    }
    & & \\ \\

    \theutterance \stepcounter{utterance}  
    & & & \multicolumn{2}{p{0.3\linewidth}}{
        \cellcolor[rgb]{0.9,0.9,0.9}{
            \makecell[{{p{\linewidth}}}]{
                \texttt{\tiny{[GM$|$GM]}}
                \texttt{game\_result = LOSE} \\
            }
        }
    }
    & & \\ \\

\end{supertabular}
}

\end{document}
